\subsection{ORB-SLAM3 Integration}
\label{sec:orb_slam3_integration}

The libVBEE is designed to be easily integrated into existing KV-SLAM systems. This section covers the integration of libVBEE into ORB-SLAM3. As stated in section \ref{sec:orb_slam3}, ORB-SLAM3 was selected due to its wide adoption, high performance, and inclusion of several SLAM extensions which can benefit from outdated map point culling. section \ref{sec:orb_slam3_integration_eval} will explore the effects of the VBEE integration on several aspects of the SLAM pipeline.

Integration begins in the \texttt{MapPoint} class, where each map point is given an instance of the VBEE class. This allows each map point to have its viewpoint-based existence probability managed separately. Each VBEE instance is initialized with the same parameter vectors $\params{\vbee}, \params{\knn}, \params{\epe}$ in the \texttt{MapPoint} constructor, ensuring that all map points share the same configuration. The VBEE instance manages the existence probability of the map point throughout its lifetime.

To update the existence probability, the map point must receive positive and negative observations. Positive observations are straightforward to retrieve: after a new frame has been processed, all the observed map points are available in the \texttt{Frame} class. Negative observations require a more complex approach, as they are not directly available. The function $\operatorname{FindMapPointsInCameraView}$, shown in Algorithm \ref{alg:map_points_in_camera_view}, is implemented in the \texttt{Map} class to retrieve all map points that are within the camera's view. This function uses the current camera pose, the camera's projection matrix, and the active \texttt{Map} instance to find all map points that would be projected into the camera's view. The function returns a vector of map points in view. Negative observations are then determined by checking whether a map point is in the camera's view but not in the frame's observed map points.

% TODO: This algorithm is poorly implemented, it should be rewritten to use the camera's projection matrix and the map point's position to determine if the point is in the camera's view.
% TODO: Consider decoupling in-view determination from negative observation classification for clarity.
\begin{singlespace}
    \begin{algorithm}[H]
        \caption[Find Map Points in Camera View]{$\operatorname{FindMapPointsInCameraView}$}
        \label{alg:map_points_in_camera_view}
        \begin{algorithmic}
            \Function{FindMapPointsInCameraView}{camera\_pose, projection\_matrix, map}
            \State camera\_view\_points $\gets$ []
            \For{map\_point in map.GetAllMapPoints()}
            \If{IsPointInCameraView(map\_point, camera\_pose, projection\_matrix)}
            \If{map\_point not in observed\_map\_points}
            \State camera\_view\_points.append(map\_point)
            \EndIf
            \EndIf
            \EndFor
            \State \Return camera\_view\_points
            \EndFunction
        \end{algorithmic}
    \end{algorithm}
\end{singlespace}

The main loop of ORB-SLAM3's Tracking thread is modified to update all map points in the camera's view after each frame is processed. The function $\operatorname{UpdateExistenceProbabilities}$, shown in Algorithm \ref{alg:update_existence_probabilities}, updates the existence probabilities of all map points in the camera's view. This function retrieves the camera's pose and projection matrix, finds all map points in the camera's view, and updates their existence probabilities based on positive and negative observations.

\begin{singlespace}
    \begin{algorithm}[H]
        \caption[Update Existence Probabilities]{$\operatorname{UpdateExistenceProbabilities}$}
        \label{alg:update_existence_probabilities}
        \begin{algorithmic}
            \Require{$frame$: Current Frame}
            \Function{UpdateExistenceProbabilities}{}
            \State{$all\_map\_points \gets \operatorname{FindMapPointsInCameraView}()$}
            \State{$seen\_points \gets$ all map points observed in $frame$}
            \State{$unseen\_points \gets$ all points in $all\_map\_points$ that are not in $seen\_points$}
            \For{$map\_point$ in $seen\_points$}
            \State{$observation \gets$ positive observation of $map\_point$ from $camera\_pose$}
            \State{$map\_point.\vbee.\update{observation}$}
            \EndFor
            \For{$map\_point$ in $unseen\_points$}
            \State{$observation \gets$ negative observation of $map\_point$ from $camera\_pose$}
            \State{$map\_point.\vbee.\update{observation}$}
            \EndFor
            \EndFunction
        \end{algorithmic}
    \end{algorithm}
\end{singlespace}

C++ implementations of these functions are provided in libVBEE for ease of integration, and can be seen in section \ref{sec:key_algorithm_impl}. With these modifications, each map point now has an existence probability estimate that is updated based on the observations received during the SLAM process. However, these probabilities are not yet used to cull outdated map points. The next section explores how culling logic is implemented into ORB-SLAM3 to make use of the existence probabilities provided by libVBEE.

\subsubsection{Culling Logic Implementation}
\label{sec:orb_slam3_culling_logic}

ORB-SLAM3 performs map point culling as part of the \texttt{LocalMapping} thread. The process looks through the list of recently added map points, and checks them against several criteria.
\begin{itemize}
    \item Another thread has marked the map point as ``bad''
    \item The ``found ratio'' drops below 0.25
    \item The map point is seen too few times in the first 3 keyframes
\end{itemize}
These criteria are effective for quickly removing false matches, but as the map point culling process only operates only recently added map points, it 
